\documentclass[11pt,a4paper]{article}

\usepackage[T1]{fontenc}
\usepackage[utf8]{inputenc}
\usepackage{lmodern}
\usepackage[margin=0.82in]{geometry}
\usepackage{microtype}
\usepackage{amsmath,amssymb}
\usepackage{booktabs,longtable,array,tabularx}
\usepackage{enumitem}
\usepackage{xcolor}
\usepackage{hyperref}
\usepackage{listings}
\usepackage{fancyhdr}

\hypersetup{colorlinks=true,linkcolor=blue!50!black,urlcolor=blue!50!black}
\setlist{nosep,leftmargin=*}
\setlength{\parindent}{0pt}
\setlength{\parskip}{0.45em}
\pagestyle{fancy}
\fancyhf{}
\lhead{LM-IDNet}
\rhead{Atomic Implementation Roadmap}
\cfoot{\thepage}
\sloppy

\newcommand{\status}{\(\square\) Not started \quad \(\square\) In progress \quad \(\square\) Blocked \quad \(\square\) Done}
\newcommand{\task}[4]{%
  \subsubsection{#1: #2}
  \textbf{Implementation.} #3

  \textbf{Verification and acceptance gate.} #4

  \textbf{Status.} \status
}
\newcommand{\artifact}[1]{\texttt{#1}}
\newcommand{\gate}[1]{\textcolor{blue!45!black}{\textbf{Phase exit gate:} #1}}

\title{LM-IDNet: Atomic Implementation Roadmap\\
\large From Research Scaffold to Reproducible Adaptive IoT Anomaly Detector}
\author{Master's Thesis Engineering Plan}
\date{\today}

\begin{document}
\maketitle

\begin{abstract}
This document converts Chapter~5 of the LM-IDNet proposal into an executable engineering plan. Every task is intentionally small, has one primary outcome, and is paired with an objective acceptance gate. The plan preserves the proposal's phases while adding software-quality, security, data-governance, numerical-reliability, and research-reproducibility controls needed for a defensible thesis artifact. A task may be marked Done only when its test passes and the named evidence is retained in version control or an immutable experiment directory.
\end{abstract}

\tableofcontents
\newpage

\section{Execution and Control Model}

\subsection{Definitions of Done}

For every task, ``Done'' means: (i) implementation and documentation are committed; (ii) the stated automated test or manual protocol passes; (iii) no pre-existing test regresses; (iv) generated evidence is stored under \artifact{reports/}, \artifact{artifacts/}, or the CI system; and (v) assumptions and deviations are recorded in an architecture decision record (ADR) or experiment manifest. A phase closes only after its exit gate passes.

\subsection{Traceability}

Use identifiers \texttt{P0--P8} for Chapter~5 phases and \texttt{F0} for foundational engineering. Each Git issue, branch, commit, test, experiment manifest, and thesis result should cite at least one task identifier. Maintain \artifact{docs/traceability.csv} with columns \texttt{requirement}, \texttt{task}, \texttt{code}, \texttt{test}, \texttt{artifact}, and \texttt{thesis\_section}.

\subsection{Mandatory Stage Gates}

\begin{longtable}{>{\bfseries}p{0.12\textwidth}p{0.80\textwidth}}
\toprule
Gate & Required evidence \\
\midrule
\endhead
F0 & Clean CI run; locked environment; buildable container; documented architecture, security model, and data policy. \\
P0 & Immutable dataset inventory and temporal split; one-command deterministic dry run; dataset summary. \\
P1 & Validated, time-indexed matrices; silent/missing-window policy; diagnostics report; processing test suite. \\
P2 & Persisted valid model with convergence history; synthetic recovery and persistence tests. \\
P3 & Pluggable backends; high-precision numerical oracle; accuracy and runtime report. \\
P4 & Leakage-free calibration; structured scores, explanations, and anomaly events; calibration report. \\
P5 & Frozen evaluation protocol; baseline comparison with uncertainty; error analysis. \\
P6 & Shadow-tested adaptive model; poisoning controls; promotion and rollback evidence. \\
P7 & Real-data forecasts and policy-controlled advisory responses; evaluation report. \\
P8 & Reproducible edge benchmark; hardened container; deployment and limitations guide. \\
\bottomrule
\end{longtable}

\section{F0: Foundational Software Engineering and Research Controls}

\task{F0.01}{Record the baseline state}{Create \artifact{docs/baseline.md} recording the current commit, supported commands, passing and failing tests, known defects, and current gaps: inactive modeling, synthetic forecasting, absent scoring, absent LM backend, and CLI inconsistency.}{Run the existing test suite from a clean checkout and paste the exact command, environment, pass/fail counts, and failure summaries into the document. A second developer must be able to reproduce the same failures.}

\task{F0.02}{Adopt a source-layout package}{Define one import strategy for \artifact{src/}; add package metadata in \artifact{pyproject.toml}; expose a console entry point such as \artifact{lm-idnet}; remove runtime \artifact{sys.path} mutation after migration.}{Install the project into a fresh virtual environment in editable mode, execute \artifact{python -c "import ..."}, invoke \artifact{lm-idnet --help}, and run tests without setting \artifact{PYTHONPATH}.}

\task{F0.03}{Lock runtime and development dependencies}{Separate runtime, test, documentation, and development dependencies; pin direct and transitive versions using a reproducible lock mechanism; declare the supported Python version.}{Create two fresh environments from the lock file and compare installed-package manifests byte-for-byte. Run the full tests in both environments.}

\task{F0.04}{Centralize typed configuration}{Define a Pydantic configuration model for paths, dates, categories, seeds, estimator settings, calibration, adaptation, and output locations. Reject unknown fields and invalid cross-field combinations.}{Unit-test missing fields, extra fields, invalid dates, duplicate categories, non-positive window sizes, invalid quantiles, and a known-valid configuration. Snapshot the normalized configuration.}

\task{F0.05}{Separate configuration from secrets}{Keep research parameters in versioned files and secrets or host-specific paths in environment variables or ignored local overrides. Add an example environment file containing no credentials.}{Run a secret scanner over the repository and verify that a clean clone operates with the documented non-secret defaults. Confirm that local overrides are ignored by Git.}

\task{F0.06}{Standardize the command-line interface}{Define stable subcommands: \artifact{preprocess}, \artifact{diagnose}, \artifact{train}, \artifact{calibrate}, \artifact{score}, \artifact{evaluate}, \artifact{adapt}, \artifact{forecast}, and \artifact{benchmark}. Make errors non-zero and machine-readable where appropriate.}{For every subcommand, test \artifact{--help}, missing arguments, invalid configuration, and a minimal success case. Assert exit codes and essential output fields.}

\task{F0.07}{Configure formatting and linting}{Adopt an automatic formatter and a fast linter with rules for imports, complexity, unsafe constructs, and unused code; store all settings in \artifact{pyproject.toml}.}{Run both tools in check-only mode over source and tests with zero violations. Add a deliberately malformed temporary fixture in a tool test to prove the check fails.}

\task{F0.08}{Introduce static type checking}{Add annotations to public interfaces and configure a type checker with strict checks for the core \artifact{algorithms}, \artifact{models}, and artifact-schema modules. Expand strict coverage incrementally to processing and evaluation.}{Type-check the strict modules with zero errors. Include type-checking in CI and demonstrate that an incompatible count-array argument is detected in a test fixture.}

\task{F0.09}{Define the testing architecture}{Configure pytest markers for \artifact{unit}, \artifact{integration}, \artifact{numerical}, \artifact{slow}, \artifact{pcap}, and \artifact{benchmark}. Establish fixture factories for count matrices, windows, configurations, and temporary artifact stores.}{Run each marker independently; ensure no test is silently uncollected and that the default fast suite excludes only explicitly marked slow or data-heavy tests.}

\task{F0.10}{Set coverage policy}{Collect branch coverage for \artifact{src/}; set an initial non-regression threshold and higher per-module thresholds for schemas, likelihood code, calibration, and decision logic.}{CI must fail when coverage drops below the configured threshold. Generate HTML and XML reports and confirm important exception branches are exercised.}

\task{F0.11}{Add property-based tests}{Use generated inputs for non-negative count matrices, category vectors, timestamps, and configuration boundaries. Define invariants such as finite outputs, non-negative persisted counts, and category-order preservation.}{Execute at least 100 generated examples per core invariant with reproducible failure seeds. Store any minimized regression example as a permanent unit test.}

\task{F0.12}{Create continuous integration}{Add GitHub Actions jobs for formatting, linting, types, unit tests, numerical tests, coverage, package build, and container build on supported Python versions. Cache only immutable dependency data.}{Open a test pull request: verify all jobs pass; then introduce a controlled lint or test failure and verify merging is blocked. Archive the successful workflow URL in \artifact{docs/baseline.md}.}

\task{F0.13}{Configure pre-commit checks}{Run formatting, linting, trailing-whitespace, large-file, merge-conflict, YAML/JSON, and secret checks before commits. Exclude only documented generated or research-data paths.}{Run hooks against all files with zero failures. Attempt to commit a synthetic secret and an oversized dummy PCAP and verify both are rejected.}

\task{F0.14}{Harden container builds}{Create a minimal non-root runtime image, pin the base image by digest, use a multi-stage build, exclude raw data and development files, and add a meaningful application health check.}{Build without network-dependent runtime steps, inspect the image user, run a read-only-root-filesystem smoke test with explicit writable mounts, and scan for critical/high vulnerabilities.}

\task{F0.15}{Make container execution reproducible}{Provide a Compose profile for development and a separate benchmark profile with fixed resource limits, read-only configuration mounts, and explicit artifact volumes.}{On a clean machine or CI runner, build and run the smoke dataset twice; compare normalized artifacts and confirm the container cannot modify configuration or raw-data mounts.}

\task{F0.16}{Adopt structured logging}{Replace global logging configuration with a library-safe logger and structured records containing run ID, stage, dataset ID, device ID, model version, and severity; never log packet payloads.}{Capture logs in tests and assert required fields, stable severity levels, and absence of payload bytes, secrets, and absolute private source paths.}

\task{F0.17}{Define exception and failure semantics}{Create domain exceptions for configuration, ingestion, validation, numerical precision, convergence, artifact compatibility, and policy rejection. Ensure commands fail closed when model integrity is uncertain.}{Unit-test every exception mapping to CLI exit code and message. Corrupt a model checksum and verify scoring stops rather than continuing with a warning.}

\task{F0.18}{Version schemas and artifacts}{Assign semantic schema versions to processed datasets, models, thresholds, events, and experiment manifests. Define compatible migrations and reject unknown major versions.}{Round-trip each schema through JSON; load the previous supported version through a migration test; reject an unsupported future major version.}

\task{F0.19}{Create immutable experiment manifests}{For every run, store run ID, UTC time, command, code commit, dirty-tree flag, configuration hash, dataset hashes, seed, environment, backend, and output checksums.}{Execute the same experiment twice and verify manifests differ only in permitted runtime fields while configuration, input, and output hashes match.}

\task{F0.20}{Establish research-data governance}{Classify raw PCAP, derived counts, models, logs, and reports. Document provenance, license/consent, retention, permitted sharing, anonymization, and deletion procedures. Treat PCAP payloads and endpoint identifiers as sensitive.}{Complete a data-protection checklist; scan a sample published artifact for payloads, MAC/IP addresses, hostnames, and absolute paths; obtain supervisor approval for the sharing classification.}

\task{F0.21}{Protect raw evidence}{Make raw PCAP inputs read-only during experiments; compute SHA-256 checksums; never rewrite them in preprocessing. Store derived data separately with provenance links.}{Hash every source file before and after preprocessing and assert equality. Attempt an in-place write under the container profile and verify permission denial.}

\task{F0.22}{Define the security threat model}{Document assets, trust boundaries, attacker capabilities, poisoning/evasion/replay risks, malicious PCAP parsing, artifact tampering, denial of service, and false-response risks. Map mitigations to tasks.}{Conduct a STRIDE-style review of the architecture. Every identified high-risk threat must have an owner, mitigation, and verification task or an explicitly accepted residual risk.}

\task{F0.23}{Write architecture decision records}{Create ADRs for category exclusivity, silent-window semantics, time-zone handling, likelihood coefficient convention, artifact formats, backend interface, split strategy, and advisory-only responses.}{Review ADRs against the proposal. Run a traceability check ensuring each listed decision has one accepted ADR and no contradictory configuration default.}

\task{F0.24}{Define reproducibility commands}{Add a task runner with single commands for bootstrap, quality checks, fast tests, full tests, preprocessing, experiment reproduction, report generation, and cleanup of disposable outputs.}{Starting from a clean clone, follow only the README commands to reproduce the smoke experiment. The process must not require undocumented manual file edits.}

\task{F0.25}{Create contribution and review policy}{Document branch naming, atomic commits, issue templates, code review, numerical-review requirements, and the rule that generated results require manifest-backed provenance.}{Use the template on one pilot change; confirm the review checklist links requirement, test, result, and documentation evidence.}

\gate{All F0 tasks pass in CI from a clean clone and the baseline, architecture, threat model, data policy, and ADRs are reviewed.}

\section{P0: Research Setup and Reproducibility}

\task{P0.01}{Inventory source captures}{Enumerate expected and present D-Link Camera5 captures with date, byte size, packet count where readable, checksum, and parser status. Record missing or duplicate files without silently skipping them.}{Compare inventory entries with the filesystem; assert unique checksums unless duplication is explained; fail the inventory command when a configured required capture is missing.}

\task{P0.02}{Verify capture integrity}{Use a PCAP validation tool or complete streaming parse to detect truncation, invalid timestamps, and corrupt records. Record recoverable warnings separately from fatal corruption.}{Parse every configured source to end-of-file. Store counts and warnings; rerun and confirm identical counts. Quarantine a deliberately truncated fixture and verify correct failure classification.}

\task{P0.03}{Confirm dataset provenance and benign assumptions}{Record dataset origin, collection conditions, device identity, known activities, and the evidence supporting ``normal'' training intervals. Explicitly label uncertainty.}{Supervisor reviews and signs off the provenance note. Any unverified-benign interval is flagged and excluded or handled by a documented robust-training experiment.}

\task{P0.04}{Resolve the October 10--11 split decision}{Reconcile proposal, ALM paper, available captures, and thesis objective; encode the accepted inclusion/exclusion choice in an ADR and configuration.}{Generate split summaries under both alternatives, document the effect on sample counts, and verify the chosen dates exactly match the approved ADR.}

\task{P0.05}{Freeze temporal partitions}{Create disjoint fit, calibration, development-test, and final-test date lists. Preserve chronology and prohibit random window mixing across partitions.}{Automated tests assert non-overlap, chronological order, capture-level grouping, and that final-test identifiers are unavailable to training and threshold code.}

\task{P0.06}{Create dataset fingerprints}{Hash the ordered file list, individual inputs, feature taxonomy, window policy, and partition definition into a dataset version identifier.}{Changing any file, category order, window length, or date assignment must change the fingerprint; repeating unchanged inputs must reproduce it.}

\task{P0.07}{Centralize deterministic seeds}{Define named seeds for simulation, baseline fitting, bootstrap intervals, and randomized algorithms. Propagate generators explicitly instead of using global state.}{Run the smoke experiment twice in separate processes and compare model parameters, injected anomalies, metrics, and normalized reports exactly or within a declared numerical tolerance.}

\task{P0.08}{Create run-directory isolation}{Write each experiment to \artifact{artifacts/runs/<run-id>/} with subdirectories for manifests, models, thresholds, events, metrics, plots, and logs. Prevent accidental overwrite.}{Start two concurrent smoke runs; verify no path collision, complete manifests, and atomic finalization. Reusing an existing run ID must fail unless an explicit disposable-run option is used.}

\task{P0.09}{Separate preprocessing from downstream commands}{Ensure \artifact{train}, \artifact{score}, and \artifact{forecast} consume versioned processed artifacts and never trigger implicit PCAP reprocessing.}{Mock the PCAP reader and invoke each downstream command; assert it is never called. Missing processed inputs must produce an actionable error.}

\task{P0.10}{Build a tiny public smoke fixture}{Create a payload-free synthetic packet/count fixture small enough for CI, with known timestamps, categories, silent windows, and expected outputs.}{Verify the fixture contains no real identifiers, its license permits distribution, and its expected preprocessing result is asserted exactly in integration tests.}

\task{P0.11}{Generate the initial dataset report}{Report dates, packets, windows, missing intervals, silent intervals, per-category totals, category taxonomy, and partition counts in machine-readable and LaTeX forms.}{Cross-check totals independently from source or processed files; assert report tables reconcile with dataset manifests and compile the LaTeX fragment.}

\task{P0.12}{Create an end-to-end smoke command}{Provide one command that validates configuration, processes the smoke fixture, diagnoses it, trains, calibrates, scores, and writes a manifest-backed report.}{Execute from a clean environment in CI; assert successful exit, expected artifact set, valid schemas, checksums, and no modification of inputs.}

\gate{The immutable partition and dataset fingerprint are approved, and the deterministic smoke pipeline succeeds twice from clean environments.}

\section{P1: Data Processing Hardening}

\task{P1.01}{Define canonical category semantics}{Choose lowercase canonical names and explicitly decide whether SSDP is exclusive of UDP or counted in both. Record category order and classification priority.}{Use packets representing TCP, ordinary UDP, SSDP/UDP:1900, ARP, and unknown protocols; assert exactly the counts required by the ADR and fixed column order.}

\task{P1.02}{Normalize category input}{Normalize parser output once at the ingestion boundary and reject category collisions after normalization. Remove mixed-case assumptions from later stages.}{Property-test mixed-case labels and whitespace; verify downstream matrices are identical and duplicate normalized configuration entries fail validation.}

\task{P1.03}{Harden timestamp parsing}{Convert all packet times to timezone-aware UTC, reject non-finite values, and preserve sufficient precision. Record original capture time-zone assumptions.}{Test epoch boundaries, fractional seconds, naive values, out-of-order packets, and daylight-saving examples; assert normalized UTC values.}

\task{P1.04}{Specify window boundaries}{Implement half-open windows \([t,t+\Delta)\), with an explicitly chosen origin and deterministic boundary behavior.}{Place packets immediately before, on, and immediately after a boundary and assert assignment to exactly one correct window for 1, 5, 10, and 30 minutes.}

\task{P1.05}{Preserve window timestamps}{Extend \artifact{WindowRecord} with device ID, start/end UTC timestamps, ordered counts, total count, missingness state, and metadata.}{Schema round-trip preserves timestamps and category order; assert \(N_t=\sum_{k=1}^{K}x_{tk}\) for every record.}

\task{P1.06}{Represent silent windows}{Materialize elapsed capture windows with zero total count and label them \artifact{observed-silent}; do not confuse them with absent capture coverage.}{Use a fixture with a deliberate quiet interval and verify all expected zero windows exist, remain ordered, and survive JSON/matrix export.}

\task{P1.07}{Represent missing coverage}{Define a separate \artifact{missing} state for gaps where observation is unavailable; prohibit automatic conversion of missing coverage to zeros.}{Use a fixture with declared capture discontinuity and assert missing windows are excluded or handled according to policy, never counted as observed silence.}

\task{P1.08}{Validate count values}{Reject booleans, negatives, fractional values, NaN, infinity, and integer overflow; use a documented integer width.}{Parameterized tests exercise every invalid type and maximum supported values. Valid arrays retain exact integer values through all exports.}

\task{P1.09}{Validate temporal streams}{Check monotonicity, duplicate windows, overlap, expected duration, device continuity, and category-schema consistency.}{Inject one defect of each type and assert a specific validation error containing dataset and window identifiers. A valid multi-day stream passes.}

\task{P1.10}{Make PCAP parsing resource-safe}{Stream packets with bounded memory, close readers on success and failure, constrain pathological packet handling, and count unsupported packets.}{Measure memory over a representative capture, inject parser exceptions, and assert file descriptors close and partial outputs are not promoted.}

\task{P1.11}{Make export atomic}{Write derived artifacts to a temporary sibling, validate them, fsync where appropriate, and rename atomically; never leave a valid-looking partial file.}{Interrupt export in a fault-injection test and verify no final artifact exists. A successful export must validate before becoming visible.}

\task{P1.12}{Export canonical JSON}{Define canonical JSON with schema version, dataset fingerprint, metadata, window records, and checksums.}{Round-trip a multi-day fixture and compare every field. Verify stable serialization and checksum across repeated exports.}

\task{P1.13}{Export modeling matrices}{Export an efficient tabular or NumPy representation containing timestamps, missingness mask, category columns, totals, device ID, and partition, without losing provenance.}{Load both canonical JSON and matrix output and assert row count, order, counts, timestamps, and metadata agree exactly.}

\task{P1.14}{Load matrices by partition}{Implement a loader that selects only declared partition identifiers and returns \(X\in\mathbb{N}^{R\times K}\) plus aligned metadata.}{Assert shape, dtype, non-negativity, category order, and exact selected dates. Attempt to request an undeclared partition and verify failure.}

\task{P1.15}{Compute descriptive statistics}{For each category compute \(\bar{x}_{\cdot k}\), sample variance \(s_k^2\), zero rate, min/max, quantiles, and total counts, with defined behavior for insufficient samples.}{Compare a hand-calculated fixture against expected values and NumPy/SciPy references within stated tolerance.}

\task{P1.16}{Compute dispersion diagnostics}{Compute the Poisson dispersion index
\[
\mathrm{DI}_k=\frac{s_k^2}{\bar{x}_{\cdot k}},
\]
returning an explicit undefined state when \(\bar{x}_{\cdot k}=0\).}{Test underdispersed, equidispersed, overdispersed, and all-zero fixtures. No NaN may enter a serialized report without an explicit status.}

\task{P1.17}{Report cross-category dependence}{Compute covariance/correlation where defined and visualize composition dependence to support the Dirichlet--Multinomial rationale.}{Compare covariance on a known matrix with an independent implementation; mark constant-category correlations undefined rather than silently zero.}

\task{P1.18}{Generate data-quality reports}{Produce JSON and LaTeX reports covering validation, missing/silent windows, unknown protocols, dispersion, totals, and partition summaries.}{Validate JSON against schema, compile LaTeX without warnings treated as errors where feasible, and reconcile every reported count with processed artifacts.}

\task{P1.19}{Add processing regression fixtures}{Freeze expected outputs for boundary, casing, SSDP classification, silent, missing, corrupt, and multi-day scenarios.}{Run regression tests on every supported Python version; any intentional output change must update an ADR and dataset schema/version.}

\task{P1.20}{Benchmark preprocessing}{Measure throughput, peak memory, unsupported-packet rate, and output size on a representative capture without including timing in correctness tests.}{Repeat at least five times after warm-up; report median and dispersion with host metadata. Confirm counts are identical in all runs.}

\gate{All configured captures produce schema-valid, time-indexed, checksum-backed matrices, and processing reports reconcile with source inventories.}

\section{P2: Dirichlet--Multinomial Modeling}

\task{P2.01}{Define estimator input contract}{Accept only two-dimensional non-negative integer matrices with \(R>1\), \(K>1\), stable category order, and a documented policy for zero-total rows and all-zero categories.}{Parameterized tests reject every invalid dimension and policy violation; a valid matrix is not mutated by training.}

\task{P2.02}{Verify the likelihood formula}{Implement the reference log probability
\begin{align}
\log p(\mathbf{x}\mid\boldsymbol{\alpha})={}&
\log\Gamma(N+1)-\sum_{k=1}^{K}\log\Gamma(x_k+1)\\
&+\log\Gamma(\alpha_0)-\log\Gamma(N+\alpha_0)
+\sum_{k=1}^{K}\left[\log\Gamma(x_k+\alpha_k)-\log\Gamma(\alpha_k)\right],
\end{align}
where \(\alpha_0=\sum_k\alpha_k\). Make inclusion of the multinomial coefficient explicit.}{Compare small cases with direct arbitrary-precision probability enumeration and verify the coefficient convention with a hand-calculated example.}

\task{P2.03}{Implement robust initialization}{Provide mean-concentration and moment-based initializers, floors for sparse categories, and deterministic fallback behavior.}{Test balanced, highly skewed, sparse, and nearly constant synthetic matrices; every returned \(\alpha_k\) must be finite and strictly positive.}

\task{P2.04}{Harden the fixed-point update}{Separate one Minka update from the iteration loop, validate numerator and denominator, guard positivity, and detect non-finite or zero updates.}{Unit-test the update against a stored high-precision example and force each guard branch with synthetic pathological inputs.}

\task{P2.05}{Define convergence criteria}{Require both parameter stability and likelihood stability, use absolute plus relative tolerances, impose maximum iterations, and distinguish converged, stalled, diverged, and exhausted outcomes.}{Construct fixtures for each outcome and assert the status, iteration count, and diagnostic reason. Never label maximum-iteration exhaustion as convergence.}

\task{P2.06}{Record optimization history}{Store iteration, \(\boldsymbol{\alpha}\), log-likelihood, absolute/relative change, update norm, and elapsed time with configurable bounded retention.}{Verify history length and final entry agree with the returned model. Serialized history must reload without loss beyond declared float precision.}

\task{P2.07}{Activate the modeling stage}{Load only the fit partition, initialize and fit the estimator, derive \(\alpha_0\), \(p_k=\alpha_k/\alpha_0\), and \(\psi=1/\alpha_0\), then persist diagnostics.}{Run the CLI on the smoke matrix; assert a valid model artifact, successful exit, correct categories, positive alpha, \(\sum_kp_k=1\), and no access to calibration/test partitions.}

\task{P2.08}{Define the model artifact}{Include schema version, model version, categories, parameters, training range, dataset/config/code hashes, backend, convergence, runtime, and integrity checksum.}{Validate against schema, round-trip exactly, reject category reordering and checksum corruption, and verify no raw packet data is embedded.}

\task{P2.09}{Recover synthetic parameters}{Sample multiple seeded datasets from known \(\boldsymbol{\alpha}^{*}\), fit them, and quantify error in mean composition and concentration.}{Across increasing sample sizes, demonstrate decreasing median estimation error or explain deviations; enforce predeclared tolerances on a stable medium-size fixture.}

\task{P2.10}{Test permutation equivariance}{Permute category columns and initialization consistently; fitted parameters should permute in the same way.}{Property-test random permutations and compare unpermuted results within numerical tolerance.}

\task{P2.11}{Test row-order invariance}{Shuffle training windows while preserving values; batch MLE results should remain unchanged up to floating-point tolerance.}{Run at least ten seeded shuffles and compare parameters and terminal likelihood with declared tolerances.}

\task{P2.12}{Assess likelihood stability}{Check for monotonic improvement where the selected fixed-point method guarantees it; otherwise define and test an accepted small numerical-decrease tolerance and stalled-run policy.}{Inspect synthetic and real histories automatically. Any decrease beyond tolerance fails training and preserves diagnostics for analysis.}

\task{P2.13}{Create model-fit diagnostics}{Report observed versus fitted mean composition, concentration, dispersion, per-window likelihood distribution, and residual summaries.}{On simulated data, verify fitted summaries approach generating values; compile JSON and LaTeX reports and reconcile model identifiers.}

\task{P2.14}{Benchmark estimator scaling}{Measure iterations, wall time, and peak memory as \(R\), \(K\), counts, and overdispersion vary.}{Use a parameterized benchmark matrix, repeat runs, and store median/quantiles with environment metadata and correctness checks.}

\gate{Training is partition-safe, produces an integrity-checked converged model, passes recovery/invariance tests, and has a reviewed fit report.}

\section{P3: Likelihood Backend Abstraction}

\task{P3.01}{Define the backend protocol}{Create a stateless \artifact{LikelihoodBackend} interface returning value, precision status, error bound when available, backend version, and timing metadata; support single-window and batch calls.}{Contract-test a fake backend and every real backend for shape, dtype, metadata, immutability, error semantics, and consistent coefficient convention.}

\task{P3.02}{Implement the standard gammaln backend}{Move SciPy-based likelihood evaluation behind the interface and support stable vectorized batch evaluation.}{Compare with the P2 arbitrary-precision oracle across small random cases and with the previous implementation on regression fixtures.}

\task{P3.03}{Build a high-precision oracle}{Use \artifact{mpmath} at configurable precision to evaluate selected log probabilities and paired log-gamma differences outside production code.}{Validate oracle precision by recomputing at doubled digits; accepted reference digits must remain unchanged. Store reference vectors and provenance.}

\task{P3.04}{Specify the LM numerical contract}{From the cited LM algorithm, document domain, requested precision, truncation/error rules, Bernoulli constants, fallback behavior, and statuses \artifact{ok}, \artifact{warning}, and \artifact{precision-impossible}.}{Review the specification line-by-line against the source paper and reproduce at least one published numerical example before implementation acceptance.}

\task{P3.05}{Implement paired log-gamma differences}{Implement the LM primitive for expressions such as \(\log\Gamma(x+a)-\log\Gamma(a)\), with no silent fallback and explicit error estimates.}{Compare a logarithmic grid of small/large \(x\), small/large \(a\), and low-overdispersion cases with the doubled-precision oracle. Enforce requested error bounds.}

\task{P3.06}{Assemble the Python LM backend}{Use the LM primitive for all paired terms in the Dirichlet--Multinomial likelihood and retain an exact/reference path for factorial terms.}{Contract-test against standard and oracle backends over randomized and adversarial grids. Assert correct precision status propagation.}

\task{P3.07}{Add optional Yu--Shaw comparison}{Implement YS only behind the same interface, using the literature-defined method and documented feasible domain; otherwise record a justified scope exclusion.}{Reproduce a cited example and compare accuracy/runtime with the oracle. Unsupported cases must return a typed status, not a plausible number.}

\task{P3.08}{Test extreme numerical regimes}{Cover zero counts, huge counts, highly imbalanced categories, \(\alpha_k\) near the lower bound, very large concentration, and cancellation-prone low-overdispersion cases.}{No accepted result may be NaN or infinite. Compare to the oracle and enforce absolute/relative/ULP tolerances appropriate to each regime.}

\task{P3.09}{Test additive batch consistency}{The sum of single-window log probabilities must agree with the batch log-likelihood under the same coefficient convention.}{Property-test random matrices and all backends; compare with a scale-aware floating-point tolerance.}

\task{P3.10}{Implement backend selection}{Select backend only through validated configuration and record backend name/version in every model, score, and manifest.}{Run identical scoring with each enabled backend and verify artifact metadata. Invalid or unavailable selections must fail before data processing.}

\task{P3.11}{Create backend benchmarks}{Benchmark accuracy, error status, latency, throughput, and memory across counts, dimensions, concentration, and batch sizes. Separate cold-start from steady-state results.}{Repeat randomized benchmark order, report median and tail latency, validate every benchmark output against the oracle subset, and archive host/toolchain details.}

\task{P3.12}{Prototype the optimized LM binding}{If justified by benchmarks, expose a C/C++ LM kernel through a narrow, memory-safe binding with explicit array ownership, bounds checking, and compiler flags.}{Use sanitizers in native tests, fuzz invalid shapes/values, compare results with Python LM/oracle, and verify no leak across repeated calls.}

\gate{Every enabled backend satisfies one interface, numerical errors are bounded against an independent oracle, and benchmark claims are reproducible.}

\section{P4: Threshold Calibration and Anomaly Scoring}

\task{P4.01}{Enforce fit/calibration isolation}{Load calibration windows through a partition-scoped API that cannot return fit or final-test rows.}{Assert disjoint identifiers and monkeypatch the fit loader to fail if calibration invokes it. Store partition fingerprints in the threshold artifact.}

\task{P4.02}{Compute raw window scores}{Return one log probability \(s_t=\log p(\mathbf{x}_t\mid\hat{\boldsymbol{\alpha}})\) and precision metadata per eligible window.}{Compare each score with the backend single-window call; preserve timestamps/order and explicitly classify missing windows.}

\task{P4.03}{Define normalized scores}{Implement a documented normalization, such as per-packet score \(s_t/\max(N_t,1)\), while treating silent windows according to ADR rather than dividing silently.}{Test totals \(N_t=0,1\), and large \(N_t\); assert finite values or explicit non-applicability and record the formula in artifacts.}

\task{P4.04}{Analyze volume confounding}{Measure the relationship between raw score and total count to decide whether raw, normalized, stratified, or paired thresholds are primary.}{Produce scatter/correlation and binned false-alarm analyses on calibration only. Record the selected policy in an ADR before final-test access.}

\task{P4.05}{Calibrate empirical quantiles}{Compute the lower-tail threshold \(\tau=Q_q(\{s_i\}_{i=1}^{m})\) with declared quantile convention and target calibration false-positive rate.}{Compare with an independent quantile implementation, test tiny/tied samples, and verify observed calibration alarm rate matches the convention within finite-sample limits.}

\task{P4.06}{Quantify threshold uncertainty}{Bootstrap calibration windows in time-aware blocks and compute a confidence interval for \(\tau\) without touching test data.}{With fixed seeds, reproduce interval endpoints; verify block resampling preserves required temporal dependence and report the effective sample size.}

\task{P4.07}{Persist threshold artifacts}{Store model checksum, score type, quantile, threshold, uncertainty, calibration range/fingerprint, sample count, backend, and schema version.}{Round-trip and checksum-test the artifact. Scoring must reject a threshold created for another model, backend convention, or category schema.}

\task{P4.08}{Implement online scoring}{Consume windows in timestamp order and emit score records without future context. Keep state bounded and deterministic.}{Replay the same stream in different chunk sizes and assert identical ordered outputs; monitor memory over a long synthetic stream.}

\task{P4.09}{Implement anomaly decisions}{Apply the strict documented rule \(s_t<\tau\) (including explicit equality behavior) and retain both score and threshold.}{Test values below, equal to, and above threshold, including floating-point neighbors. Decisions must be backend-metadata aware.}

\task{P4.10}{Calibrate severity}{Use a robust calibration scale such as IQR or median absolute deviation and define bounded severity bands from distance below threshold.}{Test degenerate calibration distributions, monotonicity, band boundaries, and invariance to irrelevant record ordering.}

\task{P4.11}{Compute category residual explanations}{For \(N>0\), calculate
\[
r_k=\frac{x_k}{N}-\frac{\alpha_k}{\alpha_0},\qquad
e_k=N\frac{\alpha_k}{\alpha_0},
\]
and rank positive and negative deviations; use a separate explanation for silence.}{Hand-check a known vector, assert residuals sum to zero within tolerance, category ordering is deterministic under ties, and silent windows never divide by zero.}

\task{P4.12}{Validate explanation faithfulness}{Distinguish residual heuristics from exact additive likelihood attribution and label them accurately in output and thesis text.}{Inject one-category increases/decreases at fixed total and verify the expected category ranks first in a predefined majority of controlled cases; document known counterexamples.}

\task{P4.13}{Define anomaly event schemas}{Create a versioned event containing device/window IDs, counts, score, threshold, decision, severity, explanation, precision status, model/threshold versions, and run ID.}{Schema-test mandatory fields, reject incompatible category order, round-trip JSON/CSV, and verify no packet payload or unnecessary endpoint identifier is present.}

\task{P4.14}{Make event output durable}{Write events atomically, support idempotent event IDs, and prevent duplicates during replay/restart.}{Interrupt and resume a stream; compare output with uninterrupted execution and assert unique event IDs and identical decisions.}

\task{P4.15}{Generate calibration reports}{Plot/table score distributions, threshold, uncertainty, target/observed alarm rate, score-volume behavior, and examples without using final-test labels.}{Regenerate from artifacts alone, compile the LaTeX fragment, and verify all figures cite model, threshold, dataset, and run identifiers.}

\gate{A model-bound threshold produces deterministic streaming decisions and schema-valid, interpretable events with documented calibration uncertainty and no data leakage.}

\section{P5: Experimental Anomaly Evaluation}

\task{P5.01}{Audit labels and attack data}{Inventory available labels, their provenance, temporal precision, attack taxonomy, and ambiguity. Do not infer ground truth from detector output.}{Cross-check labels with source documentation and sample records; report unlabeled and conflicting intervals explicitly.}

\task{P5.02}{Pre-register the evaluation protocol}{Freeze hypotheses, primary score, metrics, baselines, anomaly strengths, seeds, exclusion rules, and statistical tests before final evaluation.}{Hash and commit the protocol; evaluation code verifies the protocol hash and refuses undeclared final-test parameter tuning.}

\task{P5.03}{Define synthetic injection semantics}{Inject anomalies only into copies of held-out benign count windows; preserve timestamps, originals, labels, seed, changed categories, magnitude, and feasibility constraints.}{For every injected row, verify an untouched source row exists and a machine-readable diff matches the declared transformation exactly.}

\task{P5.04}{Implement TCP flood injections}{Increase TCP counts using absolute and multiplicative severity levels with overflow protection and fixed/variable-total variants if scientifically justified.}{Property-test non-target categories, totals, labels, determinism, and monotonic injected severity. Inspect representative cases.}

\task{P5.05}{Implement UDP burst injections}{Apply the approved UDP transformation while respecting the SSDP exclusivity ADR.}{Assert SSDP is not accidentally changed, metadata is complete, and replay with the same seed is identical.}

\task{P5.06}{Implement SSDP discovery injections}{Increase SSDP counts according to the taxonomy and category exclusivity policy.}{Verify generated counts are internally consistent with the chosen UDP/SSDP semantics and increasing severity yields non-decreasing SSDP perturbation.}

\task{P5.07}{Implement ARP scan injections}{Create temporally localized or repeated ARP count increases representing the stated abstract scan model, without claiming packet-level realism not present in counts.}{Verify duration, affected windows, magnitude, and labels exactly; document the construct-validity limitation.}

\task{P5.08}{Implement cloud-traffic suppression}{Reduce selected TCP/traffic components using bounded factors without negative counts, with an explicit rationale for the cloud proxy.}{Assert counts remain non-negative integers, reductions match factors after rounding, and zero-floor behavior is tested.}

\task{P5.09}{Implement mixed composition drift}{Shift mass among categories under fixed-total and variable-total variants, preserving integer counts and declared severity.}{Assert fixed-total variants conserve \(N\), proportions move in the intended direction, and all changes are logged.}

\task{P5.10}{Validate anomaly generators}{Review statistical magnitude and domain plausibility; report perturbations relative to benign quantiles and avoid equating synthetic anomalies with real attacks.}{Generate a validation report with before/after distributions and have the thesis supervisor approve the terminology and severity grid.}

\task{P5.11}{Implement independent Poisson baseline}{Fit one benign-rate model per category using fit data and calibrate its anomaly score on the same calibration partition.}{Recover known rates on synthetic data, test zero-rate handling, and verify no test labels affect fitting or calibration.}

\task{P5.12}{Implement Multinomial baseline}{Fit fixed composition probabilities and score counts under the same coefficient convention and partition controls.}{Compare small probabilities with a trusted reference and verify fitted probabilities sum to one.}

\task{P5.13}{Implement Gaussian baseline}{Define transformations, regularization, and singular-covariance behavior for count/proportion features.}{Recover a nonsingular synthetic Gaussian case, exercise singular data, and ensure all scores are finite or explicitly rejected.}

\task{P5.14}{Implement Isolation Forest baseline}{Use a fixed seed and a predeclared feature pipeline; tune contamination or threshold only on calibration data.}{Repeat training for deterministic scores within library guarantees, inspect feature shapes, and prove final-test labels are inaccessible.}

\task{P5.15}{Implement One-Class SVM baseline}{Scale features using fit-only statistics and select hyperparameters from a predefined calibration grid.}{Test scaler isolation, deterministic inference, and record the entire selection table rather than only the winner.}

\task{P5.16}{Standardize baseline interfaces}{Require all detectors to implement fit, calibrate, score, save/load, metadata, and explanation-capability declarations.}{Run one contract suite against every detector; compare row IDs to ensure identical evaluation populations.}

\task{P5.17}{Implement window-level metrics}{Compute confusion matrix, TPR/recall, FPR, precision, specificity, F1, balanced accuracy, ROC-AUC, and PR-AUC with explicit zero-denominator behavior.}{Compare known confusion matrices with a trusted metrics library and test single-class edge cases.}

\task{P5.18}{Implement event-level metrics}{Merge contiguous labeled anomalous windows according to a fixed rule and compute event recall, false alerts per day, and detection delay.}{Use hand-labeled timelines covering immediate, delayed, missed, duplicated, and pre-event alarms; assert exact expected results.}

\task{P5.19}{Estimate uncertainty}{Use device/day-aware or temporal-block bootstrap confidence intervals and paired resampling for detector differences.}{Simulate a known metric distribution, verify reproducibility and sensible coverage, and report resampling units and replicate count.}

\task{P5.20}{Run sensitivity analysis}{Evaluate predeclared window lengths, category sets, threshold quantiles, score normalization, and anomaly strengths without selecting on final-test performance.}{Each run must have a unique manifest; aggregate only compatible runs and show all planned settings, including failures.}

\task{P5.21}{Run ablation studies}{Remove LM backend, overdispersion modeling, normalization, explanation, and later adaptation components one at a time where meaningful.}{Verify each ablation changes only the named component by comparing normalized configuration/manifests.}

\task{P5.22}{Perform final locked evaluation}{Execute the pre-registered protocol once on the untouched final partition and store immutable outputs.}{Confirm clean tree, protocol hash, dataset fingerprint, and CI status before execution. Seal outputs with checksums and prohibit post-hoc overwriting.}

\task{P5.23}{Analyze false positives and negatives}{Sample errors by day, volume, silence/missingness, category residual, and behavioral context; distinguish model, threshold, data, and label failures.}{Use a reproducible sampling query and account for every sampled row. Record findings without changing the locked primary result.}

\task{P5.24}{Generate publication tables and figures}{Produce manifest-backed LaTeX tables for primary metrics, confidence intervals, baseline comparisons, sensitivity, and errors; generate accessible vector plots.}{Rebuild all outputs from result artifacts, compile them in a minimal LaTeX document, and cross-check selected cells against raw metrics.}

\gate{The locked protocol has run on held-out data, all baselines share partitions and metrics, uncertainty is reported, and synthetic-result claims are explicitly bounded.}

\section{P6: Adaptive Retraining}

\task{P6.01}{Define adaptation states}{Model states such as collecting, eligible, candidate-training, shadow-evaluation, promoted, rejected, and rolled-back with explicit transitions.}{State-machine tests exercise every legal transition and reject illegal transitions, retries, and duplicate promotion events.}

\task{P6.02}{Implement the safe-window predicate}{Use
\[
\mathrm{safe}(t)\Longleftrightarrow s_t\geq\tau+m
\]
with configurable margin \(m\), valid precision status, acceptable missingness, and any cooldown rules.}{Test exact boundaries, invalid precision, silent/missing windows, and anomaly cooldown periods.}

\task{P6.03}{Implement a bounded safe buffer}{Store only eligible records with IDs, scores, timestamps, model version, and integrity metadata; cap by age and count.}{Stream more than capacity, restart, and verify deterministic eviction, persistence, ordering, deduplication, and bounded storage.}

\task{P6.04}{Trigger candidate training}{Trigger only after predeclared safe-window count/time conditions and prevent overlapping updates for the same device.}{Use a simulated clock to test just-before/at/after boundaries and concurrent trigger attempts.}

\task{P6.05}{Train candidates in isolation}{Fit a candidate from the approved adaptation window without mutating the active model or threshold.}{Inject training failure and verify active artifacts remain byte-identical; successful candidates receive new immutable versions.}

\task{P6.06}{Limit parameter movement}{Constrain change in composition, concentration, or both; optionally use
\[
\boldsymbol{\alpha}_{a}=(1-\rho)\boldsymbol{\alpha}_{o}+\rho\boldsymbol{\alpha}_{n},\qquad 0<\rho\leq1.
\]}{Test \(\rho\) boundaries, positivity, deterministic output, and rejection of candidates exceeding configured drift limits.}

\task{P6.07}{Recalibrate candidate thresholds safely}{Use a candidate-specific benign validation buffer disjoint from candidate fitting data, or retain the old threshold under a documented policy.}{Assert model-threshold binding and partition disjointness; reject insufficient calibration samples.}

\task{P6.08}{Evaluate candidates in shadow mode}{Score a fixed recent validation horizon with old and candidate models without affecting alerts; compare false-alarm proxy, likelihood stability, and resource cost.}{Replay a deterministic horizon and assert promotion metrics are identical across runs and contain both model versions.}

\task{P6.09}{Define promotion policy}{Predeclare minimum improvement/non-inferiority, maximum parameter drift, convergence, numerical status, and resource limits. Require every condition.}{Boundary-test each criterion and prove a candidate failing any single condition is rejected with an auditable reason.}

\task{P6.10}{Promote atomically}{Switch active model and threshold as one transaction while retaining the previous pair and audit record.}{Fault-inject between write steps and verify readers observe either the complete old pair or complete new pair, never a mixed pair.}

\task{P6.11}{Implement automatic rollback}{Rollback when post-promotion alarm-rate, numerical, integrity, or operational guardrails fail within an observation window.}{Simulate each guardrail violation and verify atomic restoration, incident record, and cooldown against immediate re-promotion.}

\task{P6.12}{Retain model lineage}{Record parent model, safe-buffer fingerprint, training/calibration ranges, policy version, evaluation, promotion actor, and rollback reason.}{Traverse lineage from current to initial model and verify checksums and absence of cycles or missing parents.}

\task{P6.13}{Simulate gradual benign drift}{Generate or select chronologically shifting benign compositions and compare static versus adaptive false alarms and fit.}{Verify the drift generator independently, repeat across seeds, and report benefit with uncertainty and adaptation lag.}

\task{P6.14}{Simulate abrupt benign change}{Evaluate firmware-like regime change without labeling it an attack and measure recovery time and transient alarms.}{Confirm change point and magnitude; calculate recovery from a predeclared rule and compare policies.}

\task{P6.15}{Test poisoning resistance}{Inject low-and-slow, margin-crossing, burst, and mixed-category contamination into the update stream; measure parameter displacement and attack acceptance.}{Compare clean and poisoned lineage under identical seeds. Enforce a declared maximum movement or clearly report policy failure.}

\task{P6.16}{Test evasion near threshold}{Construct perturbations that remain just inside the safe margin and quantify buffer admission and later detection impact.}{Verify score placement with the active backend and document whether attacks are truly feasible at the packet level or only count-space stress tests.}

\task{P6.17}{Compare adaptation policies}{Evaluate static, full replacement, smoothing rates, buffer sizes, margins, and update frequencies under the frozen protocol.}{Use identical streams and paired uncertainty; do not choose a policy on final-test results.}

\task{P6.18}{Generate the adaptation audit report}{Report lineage, promotions/rejections, rollback events, drift response, poisoning sensitivity, compute cost, and residual risks.}{Rebuild report from audit events and artifacts only; reconcile every model version and policy decision.}

\gate{Adaptation remains off by default until shadow, atomic promotion, rollback, drift, and poisoning tests pass; the static detector remains a supported fallback.}

\section{P7: Forecasting and Response Support}

\task{P7.01}{Remove synthetic forecasting actuals}{Load chronologically held-out real window counts and preserve missing/silent semantics; prohibit sampling ``actuals'' from the fitted model.}{Mock random sampling to fail during production forecasting and verify returned observations match processed artifact IDs and counts exactly.}

\task{P7.02}{Define composition forecasts}{Forecast \(\hat{p}_{t+1,k}=\alpha_k/\alpha_0\) with the model version valid at forecast origin and prevent future leakage.}{For every forecast, verify probabilities are non-negative, sum to one, and cite an origin preceding the target window.}

\task{P7.03}{Implement moving-average total forecasts}{Predict \(N_{t+h}\) using only a predeclared trailing horizon and define behavior for missing history.}{Hand-check a short sequence and use a future-value sentinel to prove no leakage.}

\task{P7.04}{Implement seasonal total forecasts}{Create time-of-day/day-of-week benign baselines with fallback for sparse seasonal cells.}{Test calendar boundaries and unavailable cells; compare forecasts with independently grouped training data.}

\task{P7.05}{Combine total and composition}{Compute expected counts \(\hat{x}_{t+h,k}=\hat{N}_{t+h}\hat{p}_{t+h,k}\) while retaining probabilistic composition output.}{Assert category expected counts sum to predicted total within tolerance and remain non-negative.}

\task{P7.06}{Correct the Brier metric contract}{Define whether Brier score assesses multiclass probability outcomes or distribution forecasts; do not call squared distance to observed proportions a standard Brier score without qualification.}{Validate the implementation on canonical one-hot examples and document any generalized distribution-score variant separately.}

\task{P7.07}{Harden Jensen--Shannon divergence}{Normalize inputs, define zero handling, select logarithm base, and compute per-window rather than accidentally aggregating across time.}{Test identity, symmetry, finiteness, and the theoretical bound for the selected log base against trusted examples.}

\task{P7.08}{Add total-count forecast metrics}{Implement MAE, RMSE, and optionally Poisson deviance with missing-value policy.}{Compare known vectors with a trusted library or hand calculation and test all-zero targets.}

\task{P7.09}{Use rolling-origin evaluation}{Evaluate forecasts through chronological origins with fixed horizons and no refitting on future targets.}{Inspect generated origin/target pairs; automated checks require origin time less than every target time.}

\task{P7.10}{Report forecast uncertainty}{Use empirical or model-based intervals with documented assumptions and measure interval coverage.}{On seeded synthetic data, test monotonic interval endpoints and approximate expected coverage before real-data use.}

\task{P7.11}{Define advisory response taxonomy}{Create versioned response codes for monitor, operator alert, rate-limit recommendation, isolation recommendation, SSDP inspection, ARP inspection, and connectivity investigation, including reversibility, cost, and approval requirement.}{Schema-test all response records and obtain security/supervisor review that wording is advisory and avoids unsupported certainty.}

\task{P7.12}{Implement deterministic response rules}{Map severity, repetition, category residual pattern, confidence, and device criticality to recommendations using configuration rather than scattered conditionals.}{Table-test every rule, boundary, tie, and default. Same event and policy version must always yield the same recommendation.}

\task{P7.13}{Enforce advisory-only operation}{Provide no packet-blocking, firewall-changing, or device-isolation executor in the thesis prototype; mark human approval explicitly.}{Search code and container permissions for mutation paths; integration-test that response generation writes records only and performs no network-control action.}

\task{P7.14}{Audit response provenance}{Link each recommendation to anomaly event, explanation, policy version, model, and approval state.}{Trace a sample response back to checksummed source artifacts and verify policy replay reproduces it.}

\task{P7.15}{Generate forecasting and response reports}{Report composition and total metrics by horizon, baseline comparisons, uncertainty, recommendation counts, and limitations.}{Rebuild from artifacts, compile LaTeX, and cross-check sampled rows and recommendation totals.}

\gate{Forecasts use only real held-out observations and past context, metric names are statistically correct, and response output is deterministic, auditable, and strictly advisory.}

\section{P8: Edge Deployment Study}

\task{P8.01}{Define target edge profiles}{Specify representative CPU architecture, cores, memory, storage, OS, power assumptions, window arrival rate, and latency/service-level objectives.}{Review profiles against actual accessible hardware or clearly label emulation; freeze benchmark profiles before comparison.}

\task{P8.02}{Create a benchmark harness}{Measure training and per-window/batch scoring latency, throughput, peak RSS, CPU time, artifact size, startup, and optional energy, separating I/O from numerical kernels.}{Validate timers with a known-duration operation, include warm-up, repeat runs, and store raw samples and host metadata.}

\task{P8.03}{Benchmark standard Python backend}{Run the fixed workload and model under the standard gammaln backend.}{Verify output checksums against desktop reference, report median, p95/p99, dispersion, and resource limits across repeated runs.}

\task{P8.04}{Benchmark Python LM backend}{Run the identical workload under Python LM without changing model/data/configuration beyond backend.}{Manifest diff must show only the backend; numerical deviations must remain inside P3 tolerances.}

\task{P8.05}{Benchmark optimized LM backend}{Where implemented, compile reproducibly for target architecture and benchmark the same workload.}{Verify compiler/toolchain manifest, numerical equivalence, sanitizer tests, and repeated performance samples.}

\task{P8.06}{Measure bounded-memory streaming}{Replay a long stream and confirm memory stabilizes under configured buffers, logging, and adaptation settings.}{Plot RSS over time and fail if post-warm-up growth exceeds a predeclared bound after accounting for caches.}

\task{P8.07}{Measure deadline compliance}{Compare end-to-end processing and scoring latency with the configured window interval and burst arrival scenarios.}{Report deadline miss ratio and p99 latency; demonstrate backlog recovery or document infeasibility.}

\task{P8.08}{Apply least privilege}{Run as a non-root user with read-only raw/config/model mounts where possible, a dedicated writable artifact directory, dropped Linux capabilities, no-new-privileges, and no exposed port unless required.}{Inspect runtime identity/capabilities and attempt prohibited writes and privilege escalation; all must fail while scoring succeeds.}

\task{P8.09}{Secure the software supply chain}{Pin base image digest and dependencies, generate an SBOM, scan dependencies/images, document licenses, and define handling for unresolved vulnerabilities.}{CI stores SBOM and scan reports; builds fail on unwaived critical vulnerabilities and every waiver has expiry and rationale.}

\task{P8.10}{Verify offline operation}{Ensure inference requires no Internet access, telemetry, dependency download, or external service after image construction.}{Run the complete scoring smoke test with networking disabled and compare outputs with the connected run.}

\task{P8.11}{Test artifact integrity at deployment}{Verify model and threshold checksums, schema compatibility, category order, and policy version before startup; fail closed on mismatch.}{Corrupt and swap each artifact independently and verify startup rejection and a structured audit event.}

\task{P8.12}{Test graceful restart and replay}{Persist stream checkpoint and idempotent event IDs so restart neither loses nor duplicates completed-window decisions.}{Terminate at multiple processing points, restart, and compare output exactly with uninterrupted reference execution.}

\task{P8.13}{Test resource exhaustion behavior}{Apply CPU/memory/disk constraints and oversized/malformed inputs; bound queues and fail safely without corrupting artifacts.}{Run under container limits, fill the artifact volume in a controlled temporary environment, and verify explicit errors and intact previous artifacts.}

\task{P8.14}{Document deployment}{Write build, configuration, mounting, execution, health, logging, upgrade, rollback, backup, and incident procedures plus privacy warnings for PCAP data.}{A second person follows the guide on a clean host and completes offline smoke scoring and rollback without undocumented assistance.}

\task{P8.15}{Generate the edge-readiness report}{Compare backends and target profiles on correctness, latency, memory, CPU, model size, startup, and security posture; state whether near-real-time objectives are met.}{Rebuild all tables/plots from raw benchmark artifacts and manifests; independently verify selected values and confidence summaries.}

\gate{A pinned, non-root, offline-capable image passes correctness, integrity, restart, resource-limit, vulnerability, and reproducible performance tests on the declared target profile.}

\section{Cross-Phase Research and Release Controls}

\subsection{Experiment Validity Checklist}

Before any result enters the thesis, verify:
\begin{itemize}
  \item the research question and metric were declared before final-test access;
  \item fit, calibration, development, and final-test partitions are disjoint and chronological;
  \item preprocessing, category semantics, and missing/silent-window policy are identical across compared methods;
  \item baselines receive the same information and an equitable calibration budget;
  \item random seeds, negative results, exclusions, and failed runs are retained;
  \item uncertainty accompanies point estimates and the resampling unit respects temporal/device dependence;
  \item synthetic anomalies are called injections or stress tests, not evidence of detection of real attacks;
  \item plots and tables can be regenerated from immutable artifacts without hand-edited values.
\end{itemize}

\subsection{Numerical Reliability Checklist}

For every likelihood backend and estimator release:
\begin{itemize}
  \item validate \(\alpha_k>0\), \(x_k\in\mathbb{N}_0\), and compatible dimensions before computation;
  \item compare selected cases with an independent arbitrary-precision oracle;
  \item exercise cancellation-prone low-overdispersion and high-count regimes;
  \item expose, propagate, and archive precision status and error bounds;
  \item prohibit silent backend fallback or conversion of invalid output into a normal score;
  \item state whether the multinomial coefficient is included in every score and comparison.
\end{itemize}

\subsection{Release Checklist}

Tag a thesis release only when the tree is clean; CI and container scans pass; configuration, dataset, protocol, and code hashes are recorded; model and result artifacts validate; documentation compiles; licenses and data-sharing rules are satisfied; and a clean-environment reproduction has succeeded. The release notes must list limitations, unsupported environments, known numerical regimes, and any deviation from this roadmap.

\section{Optional Extensions After the Core Gates}

Theoretical extensions from Chapter~5 are deliberately deferred until P0--P5 are complete. Each extension begins with a separate hypothesis and baseline: latent traffic-topic modes; hierarchical device models; contextual parameters by time/device state; zero-inflated silence modeling; Bayesian threshold uncertainty; formal influence-bounded adaptation; or a Dirichlet--Softmax bridge. No extension may change the locked primary experiment retroactively.

\section{Final Deliverable Matrix}

\begin{longtable}{p{0.25\textwidth}p{0.45\textwidth}p{0.22\textwidth}}
\toprule
Deliverable & Minimum contents & Acceptance owner \\
\midrule
\endhead
Reproducible codebase & Typed package, stable CLI, tests, CI, locked environment, security checks & Engineer/reviewer \\
Processed dataset & Provenance, checksum, schema, timestamps, partitions, silent/missing states & Researcher/supervisor \\
Normal model & Alpha, mean, concentration, overdispersion, backend, convergence, lineage & Numerical reviewer \\
Threshold model & Model binding, calibration fingerprint, uncertainty, score convention & Researcher \\
Anomaly output & Versioned scores/events, explanations, severity, provenance & Security reviewer \\
Evaluation package & Protocol, baselines, metrics, uncertainty, injections, error analysis & Supervisor \\
Adaptive module & Safe buffer, shadow test, promotion, rollback, poisoning report & Security reviewer \\
Forecast/response module & Real actuals, validated metrics, advisory policy, audit trail & Supervisor \\
Edge package & Hardened image, SBOM, benchmark data, deployment guide & Deployment reviewer \\
Thesis evidence & Rebuildable LaTeX tables/figures tied to immutable manifests & Author/supervisor \\
\bottomrule
\end{longtable}

\section{Progress Summary}

\begin{longtable}{>{\bfseries}p{0.10\textwidth}p{0.46\textwidth}p{0.34\textwidth}}
\toprule
Phase & Outcome & Status \\
\midrule
\endhead
F0 & Engineering and research foundations & \status \\
P0 & Frozen reproducible research setup & \status \\
P1 & Validated data-processing pipeline & \status \\
P2 & Reproducible DM model training & \status \\
P3 & Reliable pluggable likelihood engine & \status \\
P4 & Calibrated online anomaly scoring & \status \\
P5 & Controlled comparative evaluation & \status \\
P6 & Guarded adaptive retraining & \status \\
P7 & Forecasting and advisory response support & \status \\
P8 & Secure edge-deployment evidence & \status \\
\bottomrule
\end{longtable}

\end{document}
